\documentclass[aspectratio=169,11pt]{beamer}
\usepackage{teaching_slides}

\title[Transformations]{ Econometrics I}
\subtitle{Lecture 2b: Transformations --- Logs, Odds, and Interpretation}
\author{Chris Conlon \\NYU Stern}
\date{Fall 2026}

% New deck (Fall 2026): why we transform y and x, tied to business settings.

\begin{document}
\maketitle


\begin{frame}{Why transform anything?}
Three distinct reasons, often confused:
\begin{enumerate}
	\item \alert{Interpretation}: you want the coefficient to answer the question
	people actually ask --- ``what's the \emph{percentage} effect?'', ``what's the
	\emph{elasticity}?''

	\medskip
	\item \alert{Functional form}: the true relationship is multiplicative or has
	diminishing returns, and a transformation makes it linear-in-parameters.

	\medskip
	\item \alert{Statistical behavior}: skewed outcomes, variance that grows with the
	level, outliers with huge leverage.
\end{enumerate}
\bigskip
Most business data (revenue, prices, firm sizes, salaries, clicks) checks all three
boxes --- which is why ``economists love logs.''
\end{frame}


\begin{frame}{The four regressions}
\begin{center}
\begin{tabular}{llll}
\hline
Model & Specification & $\beta_1$ means & Canonical business use \\
\hline
level--level & $y = \beta_0 + \beta_1 x$ & $\Delta y$ per unit $\Delta x$ & test scores on class size \\
log--level & $\ln y = \beta_0 + \beta_1 x$ & $100 \beta_1$ \% $\Delta y$ per unit $\Delta x$ & salary on years of tenure \\
level--log & $y = \beta_0 + \beta_1 \ln x$ & $\beta_1 / 100$ $\Delta y$ per 1\% $\Delta x$ & sales on ad spend \\
log--log & $\ln y = \beta_0 + \beta_1 \ln x$ & \alert{elasticity}: \% $\Delta y$ per \% $\Delta x$ & demand on price \\
\hline
\end{tabular}
\end{center}
\bigskip
\begin{itemize}
	\item The log--log slope is the \alert{price elasticity of demand} when $y$ is
	quantity and $x$ is price: $\beta_1 = -2$ means a 1\% price increase loses 2\% of
	units.
	\item Bonus: log--log coefficients don't depend on units (dollars vs.\ euros,
	ounces vs.\ liters) --- you verified this in Lecture 2.
\end{itemize}
\end{frame}


\begin{frame}{Percentage thinking: why growth is multiplicative}
\begin{itemize}
	\item Business quantities \alert{compound}: revenue grows 10\% then 10\% again,
	not \$10M then \$10M again. Asset returns, user growth, prices under inflation ---
	all multiplicative.

	\medskip
	\item Logs turn multiplicative into additive:\[
	y_t = y_0 (1+g_1)(1+g_2)\cdots \quad \Longrightarrow \quad
	\ln y_t = \ln y_0 + \textstyle\sum_t \ln(1+g_t)
	\]

	\medskip
	\item A process of many small multiplicative shocks yields a \alert{lognormal}
	outcome (Gibrat's law) --- one story for why firm sizes, incomes, and city sizes
	look lognormal-ish.

	\medskip
	\item So a linear model in $\ln y$ is the natural home for anything that grows
	by percentages.
\end{itemize}
\end{frame}


\begin{frame}{Log points vs.\ percent: when the approximation breaks}
\begin{itemize}
	\item In a log--level model, the exact effect of a one-unit change in $x$ (or a
	dummy switching on) is\[
	100 \left( e^{\beta_1} - 1 \right)\% \qquad \text{not} \qquad 100\,\beta_1 \%
	\]

	\medskip
	\item Fine for small coefficients: $\beta = 0.05$ is 5.1\%. Not fine for big ones:
	\begin{center}
	\begin{tabular}{ccc}
	$\beta$ & ``$100\beta$\%'' & exact \\
	\hline
	0.05 & 5\% & 5.1\% \\
	0.30 & 30\% & 35\% \\
	0.70 & 70\% & \alert{101\%} \\
	$-0.70$ & $-70$\% & \alert{$-50$\%} \\
	\end{tabular}
	\end{center}

	\medskip
	\item Note the asymmetry: $+0.7$ and $-0.7$ log points are $+101\%$ and $-50\%$
	--- which is also why ``up 50\% then down 50\%'' loses money.
\end{itemize}
\end{frame}


\begin{frame}{Skewness: the statistical case for logs}
\begin{itemize}
	\item Firm revenue, CEO pay, household income, order sizes, house prices:
	heavily \alert{right-skewed}. A few Walmarts and many corner stores.

	\medskip
	\item In levels, the biggest observations dominate everything: they drive the
	mean, the fit, and (via \alert{leverage}, Lecture 3) your coefficients.

	\medskip
	\item In logs, Walmart is a point like any other. Residuals look closer to normal,
	and no single firm decides your estimate.

	\medskip
	\item Related: variance usually grows with the level --- a \$10B firm's sales
	fluctuate by more \emph{dollars} than a \$1M firm's, but often by a similar
	\emph{percentage}. Logs stabilize the variance (less heteroskedasticity, same
	robust SEs anyway).

	\medskip
	\item Accounting-journal cousin: \alert{winsorizing} at the 1st/99th percentile
	--- same motivation (taming extreme observations), different tool, and it changes
	the estimand too.
\end{itemize}
\end{frame}


\begin{frame}{Log $x$ only: diminishing returns}
\begin{itemize}
	\item $y = \beta_0 + \beta_1 \ln x$: each \alert{doubling} of $x$ adds the same
	amount ($\beta_1 \ln 2$) to $y$.

	\medskip
	\item The natural shape for inputs with diminishing returns:
	\begin{itemize}
		\item Sales on \alert{advertising spend}: the tenth million of ad dollars
		buys less than the first.
		\item Output on firm size; utility on income; response rates on email volume.
	\end{itemize}

	\medskip
	\item A famous log--log version: CEO pay on firm size has an elasticity around
	$1/3$ across firms, decades, and countries --- try that one in levels and you
	learn nothing.

	\medskip
	\item Practical reading: ``moving a store from the 10th to the 90th percentile
	of foot traffic'' is often more honest than ``+1 visitor.''
\end{itemize}
\end{frame}


\begin{frame}{Dummies with log $y$: premia}
When $y$ is logged, coefficients on dummies are \alert{percentage premia}:
\begin{itemize}
	\item Union premium, gender gap (Lecture 4's workshop), CEO--MBA premium.
	\item Marketing versions: the \alert{brand premium} (branded vs.\ generic price
	gap), the \alert{verified-badge premium} on a platform, organic vs.\ conventional.
	\item Accounting versions: the Big-4 \alert{audit fee premium}; the loan-spread
	discount for firms with cleaner disclosures.
\end{itemize}
\bigskip
Remember the previous slide: report $100(e^{\beta}-1)\%$ when $\beta$ is large,
and remember these are \alert{conditional} premia --- conditional on whatever else
is in the regression.
\end{frame}


\begin{frame}{The log(0) problem}
\begin{itemize}
	\item Sales this week, patents this year, exports to a country: often
	\alert{exactly zero}, and $\ln 0 = -\infty$.

	\medskip
	\item Old fixes: $\ln(1+y)$, or the inverse hyperbolic sine
	$\mathrm{arcsinh}(y)$. Both look log-like for large $y$ and are defined at 0.

	\medskip
	\item \alert{The modern warning (Chen and Roth, QJE 2024):} with zeros, a
	``percentage effect'' is not well-defined, and estimates from log-like
	transformations \alert{depend on the units of $y$} (dollars vs.\ cents changes
	your treatment effect!). That's not a bug in the software; the estimand itself
	is incoherent.

	\medskip
	\item Better answers:
	\begin{itemize}
		\item \alert{Poisson regression} on the level of $y$ --- coefficients are
		still (semi-)elasticities, zeros are fine. (Session 13; this is why trade
		economists all switched.)
		\item Or split the question: extensive margin (any sales?) and intensive
		margin (log sales given sales $>$ 0).
	\end{itemize}
\end{itemize}
\end{frame}


\begin{frame}{Bounded outcomes: odds and log-odds}
\begin{itemize}
	\item Conversion rates, churn rates, default rates live in $[0,1]$ --- a linear
	model happily predicts $-4\%$ or $130\%$.

	\medskip
	\item Transform in two steps:\[
	p \;\longrightarrow\; \underbrace{\frac{p}{1-p}}_{\text{odds } \in (0,\infty)}
	\;\longrightarrow\; \underbrace{\ln \frac{p}{1-p}}_{\text{log-odds } \in (-\infty,\infty)}
	\]

	\medskip
	\item A 2\% conversion rate is odds of $0.0204$, log-odds of $-3.89$.
	Log-odds are unbounded and (roughly) the right scale for effects to be additive.

	\medskip
	\item Modeling log-odds as $x'\beta$ \alert{is} logistic regression --- coming in
	Session 12. Coefficients are changes in log-odds; $e^{\beta}$ is an
	\alert{odds ratio}.
\end{itemize}
\end{frame}


\begin{frame}{Rare events: relative vs.\ absolute}
\begin{itemize}
	\item When $p$ is small, odds $\approx p$, so a change in log-odds
	$\approx$ a \alert{percentage change in the probability itself}. A logit
	coefficient of $0.69$ on a rare outcome $\approx$ doubling the probability.

	\medskip
	\item This is where headlines mislead: ``doubles the risk of churn'' might mean
	$0.1\% \rightarrow 0.2\%$. Relative effects on rare events can be huge while
	absolute effects are trivial (and vice versa: a 1pp lift on a 60\% base rate is a
	modest 1.7\% relative change and an enormous amount of money).

	\medskip
	\item Discipline: for any rare-event result, report \alert{both} the relative
	effect and the baseline. ``Fraud risk doubles, from 2 in 10{,}000 to 4 in
	10{,}000.''

	\medskip
	\item Same discipline for your own regressions: a big elasticity on a tiny base
	is not a big business.
\end{itemize}
\end{frame}


\begin{frame}{Other transformations you will meet}
\begin{itemize}
	\item \alert{Standardizing} ($z$-scores): coefficients in standard-deviation
	units. Common in strategy/OB when units are survey scales with no natural meaning
	(``a 1 SD increase in management quality raises productivity by 0.2 SD'' ---
	Bloom--Van Reenen).

	\medskip
	\item \alert{Winsorizing / trimming}: standard in accounting and finance to
	control outliers in ratios (ROA, leverage). Know that it changes the estimand,
	and report the cutoffs.

	\medskip
	\item \alert{Ranks / percentiles}: robust to outliers, throws away magnitude
	(``moving from the 25th to 75th percentile of\dots'').

	\medskip
	\item \alert{Polynomials and splines}: transformations of $x$ that stay
	linear-in-parameters --- the $exper^2$ term you have already met.
\end{itemize}
\end{frame}


\begin{frame}{Checklist: to log or not to log}
Log $y$ when:
\begin{itemize}
	\item it is positive and right-skewed, grows/compounds in percentages, and the
	question is naturally ``what \% effect?''
\end{itemize}
\medskip
Think twice when:
\begin{itemize}
	\item $y$ has \alert{zeros} (use Poisson or split margins --- not $\ln(1+y)$);
	\item the decision-relevant quantity is in \alert{levels} (profit in dollars ---
	a CFO cannot spend log dollars);
	\item $y$ is a \alert{share or rate} in $[0,1]$ (think odds/logit instead).
\end{itemize}
\medskip
Always:
\begin{itemize}
	\item say in one sentence what your coefficient means, in units a manager
	understands --- if you can't, the transformation chose you, not the other way
	around.
\end{itemize}
\end{frame}


\section{Thanks!}

\end{document}
